% !PRINTED 337
\weberpage{337}

und die \gesperrt{Newton'sche N\"aherungsmethode} kommt also darauf
hinaus, dass man die Curve in dem Intervall $(\alpha, \beta)$ in erster
Ann\"aherung durch die Tangente in einem der Endpunkte ersetzt,
anstatt wie bei der Interpolationsmethode durch ihre Sehne.

% SOURCE FIGURE CROP: Bd. I registered scan 363 (printed p. 337), 600 dpi,
% crop x=180,y=900,w=2650,h=950; source-faithful raster, Fig. 21--22.
\weberfigureimage{weber-b1-fig21-22.png}{0.94}

Wenn der zweite Differentialquotient in dem Intervall $(\alpha, \beta)$
verschwindet, also die Curve einen Wendepunkt hat, so kann der
Fall eintreten, dass beide Endtangenten aus dem Intervall hinausf
\"uhren; dann ist die \gesperrt{Newton'sche} Methode also nicht an-

% SOURCE FIGURE CROP: Bd. I registered scan 363 (printed p. 337), 600 dpi,
% crop x=180,y=2280,w=2650,h=1080; source-faithful raster, Fig. 23--24.
\weberfigureimage{weber-b1-fig23-24.png}{0.94}

wendbar (Fig. 21). Es kann aber auch, wenn das Intervall $(\alpha, \beta)$
schon gen\"ugend eingeengt ist, der andere Fall eintreten, dass
beide Tangenten nach inneren Punkten des Intervalles f\"uhren
(Fig. 22). Indessen wird in diesen F\"allen immer die Regula
falsi eine bessere Ann\"aherung geben.

Wir wollen aber jetzt annehmen, dass $f'(x)$ und $f''(x)$ in
dem Intervall $(\alpha, \beta)$ nicht verschwinden. Dann \"andert die Curve
den Sinn ihrer Kr\"ummung nicht, und dann hat gewiss immer
\gesperrt{eine der beiden Endtangenten} ihren Schnittpunkt mit der
$x$-Axe im Inneren des Intervalles. Dies trifft sicher zu, wenn
die Tangente in dem Endpunkt des Intervalles genommen wird,

% Weber, Algebra. I.        22
%CONT>>
